\chapter{What is not formalized}
\label{chap:limitations}

We list the gaps between the statements above and a complete security proof.

\begin{itemize}
  \item \textbf{The real OT.} Soundness and zero knowledge are proved with the PVW OT
  (Chapter~\ref{chap:pvw}); zero knowledge is reduced to DDH by explicit distinguishers. Lean has
  no notion of efficient computation, so the efficiency of these distinguishers, like the hardness
  of DDH itself, is outside the formalization. Soundness with the Chou--Orlandi OT is formalized
  as a reduction (Theorem~\ref{thm:co-sound}): its bound contains the advantages of explicit square
  Diffie--Hellman solvers, whose efficiency, like the hardness of the problem, is outside Lean.
  \item \textbf{The hash.} Authenticity is proved with the hash as a random oracle. Outside that model,
  correlation robustness alone is not enough to prove Free XOR secure, in the sense that a fully
  black-box reduction is ruled out \cite[Sect.~3]{CKKZ12}; Choi,
  Katz, Kumaresan and Zhou show that a circular form of it suffices \cite[Sect.~4, Definition~4 and
  Theorem~1]{CKKZ12}.
  \item \textbf{Implementations.} The formalization is about the abstract garbling scheme of
  Chapter~\ref{chap:garbling}; no implementation of it is verified.
  \item \textbf{Number of queries.} The adversaries make a fixed number of queries, in a fixed
  pattern. Reducing an adversary that makes fewer queries, or interleaves them differently, to this
  form by padding with dummy queries is stated in the introduction but not formalized.
  \item \textbf{Adaptive corruptions and composition.} All statements assume static corruptions
  and concern a single execution. Security against adaptive corruptions and universal composition
  are discussed in Section~\ref{sec:corruptions} but not formalized.
  \item \textbf{Cut-and-choose.} The consistency of the garbler's own inputs across the evaluated
  copies is not modeled.
\end{itemize}
